\documentclass[11pt]{article}
\usepackage[margin=1in]{geometry}
\usepackage{amsmath,amssymb,amsfonts,mathtools,bm}
\usepackage[T1]{fontenc}
\usepackage[utf8]{inputenc}
\title{Inline and Display Mathematics}
\author{Source-backed grouped formula sample}
\date{}
\begin{document}
\maketitle
\section{Expressions}
\paragraph{Expression 1.} The following expression is taken from a source-backed formula corpus.

\begin{align*}\left({}^{c}D_{-}^{\alpha,\alpha',\beta,\beta',\gamma}f\right)(x)=\left(-1\right)^{[\operatorname{Re}(\gamma)]+1}\left(I_{-}^{-\alpha',-\alpha,-\beta',-\beta+[\operatorname{Re}(\gamma)]+1,-\gamma+[\operatorname{Re}(\gamma)]+1}f^{({[\operatorname{Re}(\gamma)]+1})}\right)(x),\end{align*}

\paragraph{Expression 2.} The following expression is taken from a source-backed formula corpus.

\begin{align*}\int_{\mathcal{V} \times \mathcal{V} } \varphi_\lambda \, d \widetilde{\pi} = \int_{\mathcal{V}} g \, d \gamma - \int_{\mathcal{S}_{1} \times \mathcal{S}_{1} } \lambda_1 c_{1} \, d \nu_{1}-\int_{\mathcal{S}_{2} \times \mathcal{S}_{2} } \lambda_2 c_{2} \, d \nu_{2} < \infty,\end{align*}

\paragraph{Expression 3.} The following expression is taken from a source-backed formula corpus.

\begin{align*}\exp \left( - i e \Bigg[ \frac{1}{\sqrt{\pi\!+\!gN}}\Big( A^{(1)}_\mu, \varepsilon_{\nu \mu} \partial_\nu \Phi^{(1)} \Big) +\sum_{I=2}^N \frac{1}{\sqrt{\pi}}\Big( A^{(I)}_\mu, \varepsilon_{\nu \mu} \partial_\nu \Phi^{(I)} \Big) \Bigg]\right) \Bigg\rangle_{\{K^{(I)}\}} \; =\end{align*}

\paragraph{Expression 4.} The following expression is taken from a source-backed formula corpus.

\begin{align*}L=i\overline{\psi}'_{0}{\gamma}^{\mu}{\partial}_{\mu}{\psi}'_{0}+i\overline{\psi}'_{0}e^{i{t}^{a}{\omega}_{a}}e^{-i{t}^{b}{\theta}_{b}}{\gamma}^{\mu}\partial_{\mu}\left (e^{i{t}^{c}{\theta}_{c}}e^{-i{t}^{l}{\omega}_{l}} \right ){\psi}'_{0}-m\overline{\psi}'_{0}{\psi}'_{0}.\end{align*}

\paragraph{Expression 5.} The following expression is taken from a source-backed formula corpus.

\begin{align*}\left(I_{-}^{\alpha,\alpha',\beta,\beta',\gamma}t^{-\rho}\right)(x)=\frac{\Gamma(-\beta+\rho)\Gamma(\alpha+\alpha'-\gamma+\rho)\Gamma(\alpha+\beta'-\gamma+\rho)}{\Gamma(\rho)\Gamma(\alpha-\beta+\rho)\Gamma(\alpha+\alpha'+\beta'-\gamma+\rho)}x^{-\alpha-\alpha'+\gamma-\rho}.\end{align*}
\end{document}
